\section{\hr{Betrieb}}


\subsection{\hr{Bedienung}}

\begin{figure}[H]
    \centering
    \includegraphics[width=6cm]{vario-display.jpg}
    \caption{LARUS Vario Display}
    \label{fig:vario-display}
\end{figure}

\tr{Das Gerät verfügt über einen Drehknopf mit zwei Ebenen und einem Drucktaster. Dem Knopf und den
beiden Drehknöpfen sind folgende Funktionen zugewiesen:}
\begin{itemize}
	\item	\tr{Drehen Sie den kleinen/oberen Drehknopf: Lautstärkeregelung}
	\item	\tr{Drehen Sie den großen/unteren Drehknopf: MacCready Wert}
	\item	\tr{Kurzer Druck auf den Drehknopf:  } \nameref{flight_menu}
	\item	\tr{Längeres Drücken des Drehknopfs: } \nameref{settings}
\end{itemize}

\tr{Das Gerät stellt drei Basisanzeigen zur Verfügung: Vario, Horizont und Geräteinformationen.
Zwischen den Anzeigen kann gewechselt werden, indem nach einem kurzen Tastendruck, im Menu der Punkt
Display angewählt wird. Die eingestellte Anzeige beibt dauerhaft erhalten. In den folgenden
Abschnitten werden die verschiedenen Anzeigemodi beschrieben. }

\subsection{Vario Display}

\tr{Die zentrale Anzeige unterscheidet zwischen dem Kreisflugmodus und dem Geradeausflugmodus. Wenn eine
Drehgeschwindigkeit von mindestens 1°/Sek. für mehr als 10 Sekunden gemessen wird, wird der
Kreisflugmodus aktiviert. Wenn die Drehgeschwindigkeit für mehr als 10 Sekunden unter diesem Wert liegt,
wird wieder in den Geradeausflugmodus zurückgeschaltet.}


\picnote
	{pictograph-blue-cloud.pdf}
	{\tr{Der Wechsel vom Variometer Modus zum Sollfahrt Modus ist unabhängig vom Wechsel der zentralen Anzeigeinformationen.}}

\tr{Weitere Informationen finden Sie im Kapitel „Umschalten von Kreisen (Vario) auf Sollfahrt“.}

\subsubsection{\hr{Kreisflugmodus}}

\begin{figure}[H]
    \centering
    \includegraphics[width=15cm]{circling-explained.pdf}
    \caption{\tr{Anzeige im Kreisflugmodus}}
    \label{fig:circling-explained}
\end{figure}

\tr{Die Anzeigen Info 1, Info 2, Info3 und die zentrale Anzeige sind davon abhängig, ob geradeaus
oder im Kreis geflogen wird. Der Inhalt dieser Anzeigen ist für beide Modus einstellbar.}

\tr{Inhalt im Kreisflugmodus:}
\begin{itemize}
	\item 	\tr{Aktuelles Steigen}
	\item 	\tr{Mittleres Steigen}
	\item 	\tr{MacCready Wert}
	\item 	\tr{Zentrale Anzeige, hier Thermikassistent}
	\item 	\tr{Info 1 Anzeige, hier Uhrzeit}
	\item 	\tr{Info 2 Anzeige, hier Windrichtung und Windstärke}
	\item 	\tr{Info 3 Anzeige, hier durchschnittliche Steigen seit Beginn des Kreisens}
	\item 	\tr{Die Farbe des Sat Symbols steht für die Datenqualität:}
	\begin{itemize}
		\item 	\tr{Grün: Verbindung zur LARUS-Sensoreinheit hergestellt. Gerät hat GPS-Fix.}
		\item	\tr{Gelb: Verbindung zur LARUS-Sensoreinheit hergestellt, Einheit hat keinen GPS-Fix}
		\item 	\tr{Rot: Keine Verbindung zur LARUS-Sensoreinheit}
	\end{itemize}
	\item 	\tr{Die Farbe des Batteriesymbols entspricht der Betriebsspannung:}
	\begin{itemize}
		\item 	\tr{Grün: Die Spannung der Batterie ist ausreichend.}
		\item	\tr{Gelb: Die Spannung der Batterie ist im kritischen Bereich.}
		\item 	\tr{Rot: Die Spannung der Batterie liegt unter dem Mindestwert.}
	\end{itemize}
\end{itemize}


\subsubsection{\hr{Geradeausflugmodus}}

\begin{figure}[H]
    \centering
    \includegraphics[width=15cm]{straight-explained.pdf}
    \caption{\tr{Anzeige im Geradeausflugmodus}}
    \label{fig:straight-explained}
\end{figure}

\tr{Inhalt im Geradeausflugmodus:}
\begin{itemize}
	\item 	\tr{Zentrale Anzeige, hier Windrichtung in bezug auf Flugzeuglängsachse}
	\item 	\tr{Mittleres Steigen}
	\item 	\tr{MacCready Wert}
	\item 	\tr{Sollfahrtgeber: Die Position des Bandes zeigt an, ob Sie zu schnell oder zu langsam
			fliegen. Positive Werte bedeuten, dass Sie zu schnell fliegen, negative Werte bedeuten,
			dass Sie zu langsam fliegen. Die Länge des Bandes gibt an, um wie viel. 1 m/s entspricht
			10 km/h.}
	\item 	\tr{Batterie Symbol ok, Sat Symbol ok}
	\item 	\tr{Info 1 Anzeige, hier Winkel Windversatz}
	\item 	\tr{Info 2 Anzeige, hier Windrichtung und Windstärke}
	\item 	\tr{Info 3 Anzeige, hier Sollfahrt}
\end{itemize}

\subsubsection{\hr{Verfügbare zentrale Anzeigen}}
\paragraph{\hr{Zentrale Anzeigen im Kreisflugmodus}}

\begin{figure}[H]
    \centering
    \includegraphics[width=5cm]{circling-single-arrow.png}
    \caption{\tr{Windanzeige mit einem Pfeil und Fahne}}
    \label{fig:circling-single-arrow}
\end{figure}

\tr{Die aktuelle Windrichtung wird durch einen zentralen Pfeil dargestellt. Die Größe des Pfeils ist
proportional zur Windgeschwindigkeit. Änderungen der Windrichtung gegenüber der mittleren
Richtung werden durch eine Windfahne angezeigt, Änderungen der Windgeschwindigkeit
gegenüber der mittelfristigen Durchschnittsgeschwindigkeit werden durch die Breite der Windfahne
dargestellt. Die Pfeilrichtung bezieht sich auf Norden.}

\begin{figure}[H]
    \centering
    \includegraphics[width=5cm]{circling-double-arrow.png}
    \caption{\tr{Windanzeige mit zwei Pfeilen}}
    \label{fig:circling-double-arrow}
\end{figure}

\tr{Die aktuelle Windrichtung und -geschwindigkeit wird durch den blauen Pfeil dargestellt,
Informationen zum mittleren Wind werden durch den grauen Pfeil im Hintergrund angezeigt. Die
Pfeilgröße hängt von der Windgeschwindigkeit ab. Die Pfeilrichtung bezieht sich auf Norden.}

\begin{figure}[H]
    \centering
    \includegraphics[width=5cm]{circling-dotted-assistant.png}
    \caption{\tr{Zentrierhilfe mit Punkten}}
    \label{fig:circling-dotted-assistant}
\end{figure}

\tr{Bedeutung der Farben der Kreispunkte:}
\begin{itemize}
	\item	\tr{Blau: Das Steigen liegt unter dem Durchschnitt}
	\item	\tr{Rot: Das Steigen liegt über dem Durchschnitt}
	\item	\tr{Gelb: Maximum des Steigens}
\end{itemize}

\tr{Der Durchmesser der Punkte ist proportional zur Steigrate.}

\begin{figure}[H]
    \centering
    \includegraphics[width=5cm]{circling-spider-assistant.png}
    \caption{\tr{Zentrierhilfe mit Spinnennetz}}
    \label{fig:circling-spider-assistant}
\end{figure}

\tr{Bedeutung der Farben der Kreispunkte:}
\begin{itemize}
	\item	\tr{Blau: Das Steigen liegt unter dem Durchschnitt}
	\item	\tr{Rot: Das Steigen liegt über dem Durchschnitt}
	\item	\tr{Gelb: Maximum des Steigens}
\end{itemize}

\tr{Der Durchmesser des Segments ist proportional zur Steiggeschwindigkeit.}

\paragraph{\hr{Zentrale Anzeigen im Geradeausflugmodus}}

\begin{figure}[H]
    \centering
    \includegraphics[width=5cm]{straight-single-arrow.png}
    \caption{\tr{Windanzeige im Geradeausflug mit einem Pfeil und Fahne}}
    \label{fig:straight-single-arrow}
\end{figure}

\tr{Die aktuelle Windrichtung wird durch einen zentralen Pfeil dargestellt. Die Größe des Pfeils ist
proportional zur Windgeschwindigkeit. Änderungen der Windrichtung gegenüber der mittleren Richtung
werden durch eine Windfahne angezeigt, Änderungen der Windgeschwindigkeit gegenüber der
mittelfristigen Durchschnittsgeschwindigkeit werden durch die Breite der Windfahne dargestellt. Das
Flugzeugsymbol deuted an, dass sich die Darstellng auf die Flugrichtung bezieht.}

\begin{figure}[H]
    \centering
    \includegraphics[width=5cm]{straight-double-arrow.png}
    \caption{\tr{Windanzeige im Geradeausflug mit zwei Pfeilen}}
    \label{fig:straight-double-arrow}
\end{figure}

\tr{Die aktuelle Windrichtung und -geschwindigkeit wird durch den blauen Pfeil dargestellt,
Informationen zum mittleren Wind werden durch den grauen Pfeil im Hintergrund angezeigt. Die
Größe der Pfeile hängt von der Windgeschwindigkeit ab. Die Pfeilrichtung bezieht sich auf die
Längsachse des Flugzeugs. Das Flugzeugsymbol deuted an, dass sich die Darstellng auf die
Flugrichtung bezieht.}

\subsection{\hr{Künstlicher Horizont}}

\picnote
	{pictograph-yellow-warning.pdf} {\tr{Dies ist kein amtlich zugelasssener künstlicher Horizont.
	Dashalb darf diese Anzeige nicht dazu benutz werden in Wolken oder anderweitig außerhalb von VFR
	Bedingeungen zu fliegen.}}

\begin{figure}[H]
    \centering
    \includegraphics[width=5cm]{horizon.png}
    \caption{\tr{Künstlicher Horizont}}
    \label{fig:horizont}
\end{figure}

\tr{Die Anzeige des künstlichen Horizonts beinhaltet folgende Informationen:}
\begin{itemize}
	\item \tr{Die blaue Fläche stellt den Himmel dar.}
	\item \tr{Die Grenze zur brauenen Fläche entspricht dem Horizont.}
	\item \tr{Die kreisförmige Skala oben zeigt in 15° Schritten die aktuelle
			Querneigung des Flugzeugs (hier etwa 30°, rote Pfeilspitze).}
	\item \tr{Die Skala parallel zum Horizont steht für den Steig- / Sinkwinkel in 10° Schritten
			(hier 0°).}
	\item \tr{Im unteren Bereich ist eine Libelle zu sehen, die ein eventuelles Schieben sichtbar
			macht.}
\end{itemize}

\tr{Das Flugzeug befindet sich derzeit in einer sauberen Rechtskurve mit einer Querneigung von 30°.}

\picnote
	{pictograph-yellow-warning.pdf} {\tr{In manchen Wettbewerben ist die Anzeige des künstlichen
	Horizontes verboten. Deshalb kann in der LARUS Sensoreinheit die Ausgabe der
	Horizontinformationen blockiert werden. In diesem Fall wird statt des Horizontes eine Warnung
	ausgegeben.}}

\subsection{\hr{Geräteinformationen}}

\tr{In dem Anzeigemodus Geräteinformationen werden Dateilinformationen zum LARUS Vario Display und
zur LARUS Sensoreinheit angezeigt. Diese Anzeige kann hilfreich sein, um Fehleranalysen
durchzuführen oder spezielle Informationen abzurufen. Beispielsweise können hier sämtliche Zustände
der Ein- und Ausgänge eingesehen werden. Auch sind Fehlerzustände der LARUS Sensoreinheit
erkennbar.}

\begin{figure}[H]
    \centering
    \includegraphics[width=5cm]{device-info.png}
    \caption{\tr{Geräteinformationen}}
    \label{fig:device-info}
\end{figure}

\subsection{Flight Menu}
\label{flight_menu}

\tr{Das }Flight Menu\tr{ stellt Einstellmöglichkeiten bereit, die vor dem Flug oder während des
Fluges benötigt werden. Das Menu erreicht man durch kurzes Drücken des Bedienknopfes.} 

\subsubsection{Water Ballast}
\tr{Hier wird die Menge des getankten Wasserballastes vorgegeben. Die Polare wird entsprechend angepasst
und bei der Berechnung der Sollfahrt berücksichtigt. Beim Ablassen des Wassers während des Fluges
kann dieser Wert manuell korrigiert oder durch einen Schalter am Ventil und entsprechender
Konfiguration automatisch reduziert werden.

Der eingestellte Wert wird zu einem angeschlossenen Navigationsrechner synchronisiert.}

\subsubsection{Bugs}
\tr{Insekten auf Tragflächen und Rumpf reduzieren die Gleitleistungen des Segelflugzeuges. Diese
Veränderungen der Leistung eines Segelflugzeuges können näherungsweise mit dieser Einstellung
berücksichtigt werden. Es sind Vorgaben von 0 bis 50 Prozent möglich. Bei einer Einstellung von
50 Prozent verdoppelt sich die Sinkrate bei gegebener Geschwindigkeit. 

Der Algorithmus arbeit exakt identisch zu XCSoar / OpenSoar. Der eingestellte Wert wird zu einem
angeschlossenen Navigationsrechner synchronisiert.}

\subsubsection{Pilot Weight}
\tr{Das Gewicht des Piloten wird bei der Berechnung der Flugzeug Polare berücksichtigt. Bei Doppelsitzern
ist hier das Gewicht der Summe beider Piloten vorzugeben

Der eingestellte Wert wird zu einem angeschlossenen Navigationsrechner synchronisiert.}

\subsubsection{Display}
\tr{Der Menupunkt Display ermöglicht die Auswahl der ständigen Anzeige des Gerätes. Hier kann zwischen
Vario, Horizont und Geräteinformationen gewechselt werden.}

